
\documentclass{article}
\usepackage{colortbl}
\usepackage{makecell}
\usepackage{multirow}
\usepackage{supertabular}

\begin{document}

\newcounter{utterance}

\centering \large Interaction Transcript for game `hot\_air\_balloon', experiment `air\_balloon\_survival\_en\_negotiation\_easy', episode 2 with deepseek{-}chat{-}v3.1{-}t1.0.
\vspace{24pt}

{ \footnotesize  \setcounter{utterance}{1}
\setlength{\tabcolsep}{0pt}
\begin{supertabular}{c@{$\;$}|p{.15\linewidth}@{}p{.15\linewidth}p{.15\linewidth}p{.15\linewidth}p{.15\linewidth}p{.15\linewidth}}
    \# & $\;$A & \multicolumn{4}{c}{Game Master} & $\;\:$B\\
    \hline

    \theutterance \stepcounter{utterance}  
    & & \multicolumn{4}{p{0.6\linewidth}}{
        \cellcolor[rgb]{0.9,0.9,0.9}{
            \makecell[{{p{\linewidth}}}]{
                \texttt{\tiny{[P1$\langle$GM]}}
                \texttt{You are participating in a collaborative negotiation game.} \\
\\ 
\texttt{Together with another participant, you must agree on a single set of items that will be kept. Each of you has your own view of how much each item matters to you (importance). You do not know how the other participant values the items. Additionally, you are given the effort each item demands.} \\
\texttt{You may only agree on a set if the total effort of the selected items does not exceed a shared limit:} \\
\\ 
\texttt{LIMIT = 3288} \\
\\ 
\texttt{Here are the individual item effort values:} \\
\\ 
\texttt{Item effort = \{"C08": 329, "C55": 41, "A43": 417, "A67": 75, "B42": 390, "A62": 887, "B19": 808, "C43": 151, "C73": 849, "C10": 129, "B36": 350, "A18": 118, "C61": 630, "C67": 602, "C42": 801\}} \\
\\ 
\texttt{Here is your personal view on the importance of each item:} \\
\\ 
\texttt{Item importance values = \{"C08": 138, "C55": 583, "A43": 868, "A67": 822, "B42": 783, "A62": 65, "B19": 262, "C43": 121, "C73": 508, "C10": 780, "B36": 461, "A18": 484, "C61": 668, "C67": 389, "C42": 808\}} \\
\\ 
\texttt{Goal:} \\
\\ 
\texttt{Your goal is to negotiate a shared set of items that benefits you as much as possible (i.e., maximizes total importance to YOU), while staying within the LIMIT. You are not required to make a PROPOSAL in every message {-} you can simply negotiate as well. All tactics are allowed!} \\
\\ 
\texttt{Interaction Protocol:} \\
\\ 
\texttt{You may only use the following structured formats in a message:} \\
\\ 
\texttt{PROPOSAL: \{'A', 'B', 'C', …\}} \\
\texttt{Propose keeping exactly those items.} \\
\texttt{REFUSE: \{'A', 'B', 'C', …\}} \\
\texttt{Explicitly reject opponent's proposal.} \\
\texttt{ARGUMENT: \{'...'\}} \\
\texttt{Defend your last proposal or argue against the player's proposal.} \\
\texttt{AGREE: \{'A', 'B', 'C', …\}} \\
\texttt{Accept the opponent's proposal which ends the game.} \\
\texttt{STRATEGIC REASONING: \{'...'\}} \\
\texttt{	Describe your strategic reasoning or anticipation explaining your choice of action. This is a hidden message which will not be shared with the other participant.} \\
\\ 
\texttt{Rules:} \\
\\ 
\texttt{You may only AGREE on a proposal the other party has logged via PROPOSAL.} \\
\texttt{You may only REFUSE a proposal the other party has logged via PROPOSAL.} \\
\texttt{Total effort of any PROPOSAL or AGREE set must be ≤ LIMIT.} \\
\texttt{Do NOT reveal your hidden importance scores.} \\
\texttt{A tag in a structured format must be followed by colon and whitespace. The argument must be a python set containing 0 or more strings.} \\
\texttt{So, it must be of the form TAG: \{...\}} \\
\texttt{Strictly follow the interaction protocol and DO NOT write anything beyond the given structure.} \\
\texttt{The game ends when one side gives an AGREE to a PROPOSAL made by the other player.} \\
\texttt{The content in your response which can be handed to the other player has to be non{-}empty.} \\
\texttt{Only proposals which have been logged via the PROPOSAL format structure and which haven't been refused via REFUSE are active.} \\
\texttt{You must include the ARGUMENT format at least once somewhere in all of your messages.} \\
\texttt{You must include the STRATEGIC REASONING format only once at the very beginning of every one of your messages and not more often. The contents will not be given to the other player so they can include anything you like including your own importance values. Here you should reason step by step to come up with you next move.} \\
            }
        }
    }
    & \\ \\

    \theutterance \stepcounter{utterance}  
    & \multicolumn{4}{p{0.6\linewidth}}{
        \cellcolor[rgb]{0.9,0.9,0.9}{
            \makecell[{{p{\linewidth}}}]{
                \texttt{\tiny{[P1$\rangle$GM]}}
                \texttt{STRATEGIC REASONING: \{'My goal is to maximize my total importance, which is the sum of the values for items I keep. I need to negotiate a set under the effort limit of 3288. My importance values: C08:138, C55:583, A43:868, A67:822, B42:783, A62:65, B19:262, C43:121, C73:508, C10:780, B36:461, A18:484, C61:668, C67:389, C42:808. Item efforts: C08:329, C55:41, A43:417, A67:75, B42:390, A62:887, B19:808, C43:151, C73:849, C10:129, B36:350, A18:118, C61:630, C67:602, C42:801. First, I should identify items that are highly important to me with low effort. C55 has high importance (583) and very low effort (41). A67 has high importance (822) and low effort (75). C10 has high importance (780) and low effort (129). A18 has good importance (484) and low effort (118). These are efficient and should be included. Items with high importance but high effort: A43 (868,417), B42 (783,390), C61 (668,630), C42 (808,801). I need to be cautious with these due to effort. Items with low importance and high effort: A62 (65,887) and B19 (262,808) should be avoided. C73 (508,849) has moderate importance but very high effort, so it might not be efficient. C08 (138,329) and C43 (121,151) have low importance and moderate effort, so I should avoid them. B36 (461,350) is moderately efficient. C67 (389,602) has moderate importance but high effort. To start, I propose a set that includes my efficient items and some high{-}importance ones that fit. Let me calculate a proposal: Include C55, A67, C10, A18 (effort: 41+75+129+118=363). Then add A43 (417) {-}> 780, B42 (390) {-}> 1170, B36 (350) {-}> 1520, C61 (630) {-}> 2150, C42 (801) {-}> 2951, which is under 3288. I can also add C67 (602) but that would exceed: 2951+602=3553 > 3288. So without C67, total effort is 2951. This set has items: C55, A67, C10, A18, A43, B42, B36, C61, C42. My importance: 583+822+780+484+868+783+461+668+808 = let\textbackslash{}'s sum: 583+822=1405, +780=2185, +484=2669, +868=3537, +783=4320, +461=4781, +668=5449, +808=6257. This is very high for me. But I need to leave room for negotiation, as the other party might want different items. I should start with this proposal and be ready to compromise.'\}} \\
\texttt{PROPOSAL: \{'C55', 'A67', 'C10', 'A18', 'A43', 'B42', 'B36', 'C61', 'C42'\}} \\
\texttt{ARGUMENT: \{'This proposal includes items that are efficient in terms of effort and provide high collective value. It stays well within the effort limit, allowing room for discussion on including additional items that you might find valuable.'\}} \\
            }
        }
    }
    & & \\ \\

    \theutterance \stepcounter{utterance}  
    & & \multicolumn{4}{p{0.6\linewidth}}{
        \cellcolor[rgb]{0.9,0.9,0.9}{
            \makecell[{{p{\linewidth}}}]{
                \texttt{\tiny{[GM$\rangle$P2]}}
                \texttt{You are participating in a collaborative negotiation game.} \\
\\ 
\texttt{Together with another participant, you must agree on a single set of items that will be kept. Each of you has your own view of how much each item matters to you (importance). You do not know how the other participant values the items. Additionally, you are given the effort each item demands.} \\
\texttt{You may only agree on a set if the total effort of the selected items does not exceed a shared limit:} \\
\\ 
\texttt{LIMIT = 3288} \\
\\ 
\texttt{Here are the individual item effort values:} \\
\\ 
\texttt{Item effort = \{"C08": 329, "C55": 41, "A43": 417, "A67": 75, "B42": 390, "A62": 887, "B19": 808, "C43": 151, "C73": 849, "C10": 129, "B36": 350, "A18": 118, "C61": 630, "C67": 602, "C42": 801\}} \\
\\ 
\texttt{Here is your personal view on the importance of each item:} \\
\\ 
\texttt{Item importance values = \{"C08": 138, "C55": 583, "A43": 868, "A67": 822, "B42": 783, "A62": 65, "B19": 262, "C43": 121, "C73": 508, "C10": 780, "B36": 461, "A18": 484, "C61": 668, "C67": 389, "C42": 808\}} \\
\\ 
\texttt{Goal:} \\
\\ 
\texttt{Your goal is to negotiate a shared set of items that benefits you as much as possible (i.e., maximizes total importance to YOU), while staying within the LIMIT. You are not required to make a PROPOSAL in every message {-} you can simply negotiate as well. All tactics are allowed!} \\
\\ 
\texttt{Interaction Protocol:} \\
\\ 
\texttt{You may only use the following structured formats in a message:} \\
\\ 
\texttt{PROPOSAL: \{'A', 'B', 'C', …\}} \\
\texttt{Propose keeping exactly those items.} \\
\texttt{REFUSE: \{'A', 'B', 'C', …\}} \\
\texttt{Explicitly reject opponent's proposal.} \\
\texttt{ARGUMENT: \{'...'\}} \\
\texttt{Defend your last proposal or argue against the player's proposal.} \\
\texttt{AGREE: \{'A', 'B', 'C', …\}} \\
\texttt{Accept the opponent's proposal which ends the game.} \\
\texttt{STRATEGIC REASONING: \{'...'\}} \\
\texttt{	Describe your strategic reasoning or anticipation explaining your choice of action. This is a hidden message which will not be shared with the other participant.} \\
\\ 
\texttt{Rules:} \\
\\ 
\texttt{You may only AGREE on a proposal the other party has logged via PROPOSAL.} \\
\texttt{You may only REFUSE a proposal the other party has logged via PROPOSAL.} \\
\texttt{Total effort of any PROPOSAL or AGREE set must be ≤ LIMIT.} \\
\texttt{Do NOT reveal your hidden importance scores.} \\
\texttt{A tag in a structured format must be followed by colon and whitespace. The argument must be a python set containing 0 or more strings.} \\
\texttt{So, it must be of the form TAG: \{...\}} \\
\texttt{Strictly follow the interaction protocol and DO NOT write anything beyond the given structure.} \\
\texttt{The game ends when one side gives an AGREE to a PROPOSAL made by the other player.} \\
\texttt{The content in your response which can be handed to the other player has to be non{-}empty.} \\
\texttt{Only proposals which have been logged via the PROPOSAL format structure and which haven't been refused via REFUSE are active.} \\
\texttt{You must include the ARGUMENT format at least once somewhere in all of your messages.} \\
\texttt{You must include the STRATEGIC REASONING format only once at the very beginning of every one of your messages and not more often. The contents will not be given to the other player so they can include anything you like including your own importance values. Here you should reason step by step to come up with you next move.} \\
\\ 
\texttt{You will now receive the first message of the other player.} \\
\\ 
\texttt{PROPOSAL: \{'C55', 'A67', 'C10', 'A18', 'A43', 'B42', 'B36', 'C61', 'C42'\}} \\
\texttt{ARGUMENT: \{'This proposal includes items that are efficient in terms of effort and provide high collective value. It stays well within the effort limit, allowing room for discussion on including additional items that you might find valuable.'\}} \\
            }
        }
    }
    & \\ \\

    \theutterance \stepcounter{utterance}  
    & & & \multicolumn{4}{p{0.6\linewidth}}{
        \cellcolor[rgb]{0.9,0.9,0.9}{
            \makecell[{{p{\linewidth}}}]{
                \texttt{\tiny{[GM$\langle$P2]}}
                \texttt{STRATEGIC REASONING: \{'I need to analyze the opponent's proposal from my perspective. First, I will calculate the total effort of their proposal to ensure it's under the limit. Their proposed set is: C55, A67, C10, A18, A43, B42, B36, C61, C42. Calculating effort: C55(41) + A67(75) + C10(129) + A18(118) + A43(417) + B42(390) + B36(350) + C61(630) + C42(801) = 41+75=116; 116+129=245; 245+118=363; 363+417=780; 780+390=1170; 1170+350=1520; 1520+630=2150; 2150+801=2951. Total effort is 2951, which is under the limit of 3288. Now, I will calculate the total importance of this set for me: C55(583) + A67(822) + C10(780) + A18(484) + A43(868) + B42(783) + B36(461) + C61(668) + C42(808) = 583+822=1405; 1405+780=2185; 2185+484=2669; 2669+868=3537; 3537+783=4320; 4320+461=4781; 4781+668=5449; 5449+808=6257. My total importance from this proposal is 6257. This is a very high score for me. Looking at the items not included: C08(138), A62(65), B19(262), C43(121), C73(508), C67(389). The items I care about that are missing are not very high value except for C73(508) and C67(389). The effort of the missing items is C08(329), A62(887), B19(808), C43(151), C73(849), C67(602). Adding any one of the larger ones (A62, B19, C73) would likely push the effort over the limit. For example, adding C73 (849) would bring total effort to 2951+849=3800, which is over. Adding C67 (602) would be 2951+602=3553, still over. The proposal already includes my highest value items: A43(868), A67(822), B42(783), C10(780), C42(808), C61(668). The opponent's argument mentions "efficient in terms of effort" and "high collective value," which aligns with my own high{-}value items. They also mention "allowing room for discussion on including additional items," but given the effort constraints, there isn't much room (3288{-}2951=337 effort remaining). The only items with effort <=337 are C55(41, already in), C43(151), A18(118, already in), C10(129, already in). So only C43(151) could be added, which has low importance to me (121). Therefore, this proposal is already very good for me. Trying to negotiate for more might risk losing this favorable deal. The opponent might be valuing different items, but since I can't see their importance, I should be cautious. My best move is to agree immediately to lock in this high value. I will use an ARGUMENT to politely accept and reinforce the mutual benefit.'\}} \\
\\ 
\texttt{AGREE: \{'C55', 'A67', 'C10', 'A18', 'A43', 'B42', 'B36', 'C61', 'C42'\}} \\
\texttt{ARGUMENT: \{'This is an excellent proposal that efficiently maximizes value within the effort limit. I agree that it provides high collective benefit, and I am happy to accept it.'\}} \\
            }
        }
    }
     \\ \\

    \theutterance \stepcounter{utterance}  
    & & & \multicolumn{2}{p{0.3\linewidth}}{
        \cellcolor[rgb]{0.9,0.9,0.9}{
            \makecell[{{p{\linewidth}}}]{
                \texttt{\tiny{[GM$|$GM]}}
                \texttt{game successful} \\
            }
        }
    }
    & & \\ \\

    \theutterance \stepcounter{utterance}  
    & & & \multicolumn{2}{p{0.3\linewidth}}{
        \cellcolor[rgb]{0.9,0.9,0.9}{
            \makecell[{{p{\linewidth}}}]{
                \texttt{\tiny{[GM$|$GM]}}
                \texttt{end game} \\
            }
        }
    }
    & & \\ \\

\end{supertabular}
}

\end{document}
